\chapter{Extraits de code et accès au dépôt}
\label{annexe:code}

\section{Accès au code}

Le code source du lab et du pipeline HOMO-CI est dans le dépôt~:

\begin{itemize}
  \item \texttt{soc/infra/}~: scripts de déploiement Azure (Bash et
    PowerShell) pour les 3 VMs.
  \item \texttt{soc/corpnet/}~: code PHP de l'application
    CorpNet (portail interne, volontairement vulnérable).
  \item \texttt{soc/internal-apps/}~: code Node.js des
    applications MaFormation et MaCandidature, plus la
    configuration Docker Compose.
  \item \texttt{soc/pipeline/}~: code Python du pipeline
    HOMO-CI.
\end{itemize}

\section{Démarrage rapide du pipeline}

\begin{lstlisting}[language=bash,caption={Installation et premier
cycle.}]
cd pipeline
python3.11 -m venv .venv
source .venv/bin/activate
pip install -r requirements.txt

export AZURE_SUBSCRIPTION_ID=...
export RESOURCE_GROUP=RG-PFE-SOC

# Cycle complet
make pipeline-run

# Sortie : build/dossier-AAAAMMJJ-HHMMSS.pdf
\end{lstlisting}

\section{Le modèle \texttt{Evidence}}

\begin{lstlisting}[language=Python,caption={Modèle Pydantic
\texttt{Evidence} (extrait).}]
class Evidence(BaseModel):
    model_config = ConfigDict(frozen=True, extra="forbid")

    id: UUID = Field(default_factory=uuid4)
    ts: datetime = Field(
        default_factory=lambda: datetime.now(timezone.utc)
    )
    source: SourceType
    category: Category
    ref: str = Field(..., min_length=1, max_length=255)
    scope: Scope = Field(default_factory=Scope)
    payload: dict[str, Any]
    payload_hash: str = Field(..., pattern=r"^[a-f0-9]{64}$")
    weight: Weight = Field(default=Weight.AUTO)
    materiality: Materiality = Field(default=Materiality.MEDIUM)

    @classmethod
    def from_payload(cls, source, category, ref, payload, ...):
        canonical = json.dumps(
            payload, sort_keys=True, separators=(",", ":")
        )
        payload_hash = hashlib.sha256(
            canonical.encode("utf-8")
        ).hexdigest()
        return cls(
            source=source, category=category, ref=ref,
            payload=payload, payload_hash=payload_hash, ...
        )
\end{lstlisting}

\section{Le diff}

\begin{lstlisting}[language=Python,caption={Moteur de diff (extrait).}]
def diff(previous, current) -> list[Change]:
    changes = []
    prev_keys = set(previous.keys())
    curr_keys = set(current.keys())

    for key in curr_keys - prev_keys:
        ev = current[key]
        changes.append(
            Change(type=ChangeType.CREATED,
                   ref=ev.ref, category=ev.category,
                   after=ev)
        )
    for key in prev_keys - curr_keys:
        ev = previous[key]
        changes.append(
            Change(type=ChangeType.DELETED,
                   ref=ev.ref, category=ev.category,
                   before=ev)
        )
    for key in prev_keys & curr_keys:
        if previous[key].payload_hash != current[key].payload_hash:
            changes.append(
                Change(type=ChangeType.UPDATED,
                       ref=current[key].ref,
                       category=current[key].category,
                       before=previous[key], after=current[key])
            )
    return changes
\end{lstlisting}

\section{Workflow GitHub Actions (cycle nightly)}

\begin{lstlisting}[language=bash,caption={\texttt{.github/workflows/nightly.yml}
(extrait).}]
name: HOMO-CI nightly
on:
  schedule:
    - cron: '0 3 * * *'   # 03:00 UTC tous les jours
  workflow_dispatch:

jobs:
  run-pipeline:
    runs-on: ubuntu-latest
    steps:
      - uses: actions/checkout@v4
      - uses: actions/setup-python@v5
        with:
          python-version: '3.11'
      - run: pip install -r pipeline/requirements.txt
      - name: Run pipeline
        env:
          AZURE_SUBSCRIPTION_ID: ${{ secrets.AZURE_SUB_ID }}
        run: |
          cd pipeline
          homo-ci collect
          homo-ci diff
          homo-ci render
          homo-ci publish
\end{lstlisting}

\section{Dépendances principales}

\begin{itemize}
  \item \texttt{pydantic >= 2.5}~: validation des modèles
    \texttt{Evidence}, \texttt{Change}, \texttt{DossierSnapshot}.
  \item \texttt{azure-identity}, \texttt{azure-mgmt-network},
    \texttt{azure-mgmt-resource}~: SDK Azure officiel.
  \item \texttt{jinja2 >= 3.1}~: templating Markdown.
  \item \texttt{pypandoc} et \texttt{weasyprint}~: conversion
    Markdown vers PDF.
  \item \texttt{click} et \texttt{rich}~: CLI et affichage.
\end{itemize}
